% !TEX root = ../../../main.tex
% Back-right column. Patterns from all 32 questions of HW01--HW04.
\section{Common proof moves}
\label{sec:appendix:formula-sheet-proof-moves}
\formulagroup{Sets, bijections, and enumeration}
Chase an arbitrary element in both directions; swap $\forall,\exists$
when negating. For a bijection, give an inverse or show disjoint image
pieces cover the target. For countable unions, read diagonals of an
array of enumerations and skip repeats.

\formulagroup{Irrationality and transcendence}
Assume a lowest-terms fraction. For $p^{1/k}$, $p$ prime, $k\ge2$,
raise to the $k$th power and force $p$ to divide both numerator and
denominator. For powers of transcendental $\tau$, clear denominators
in a nontrivial rational relation to get a nonzero polynomial vanishing
at $\tau$. Count algebraic numbers using bounded degree and coefficients;
then $\R=\mathcal A\cup\mathcal T$ forces $\mathcal T$ uncountable.
Under CH, $\aleph_0<|\mathcal T|\le|\R|$ implies equality with $|\R|$.

\formulagroup{Best bounds, epsilon choices, and cuts}
Prove a bound first, then give witnesses within every $\varepsilon$.
If $u>v$, take $\varepsilon=(u-v)/2$.
If $|x_n-L|\le C/n^\alpha$, choose $N>(C/\varepsilon)^{1/\alpha}$.
For two errors budget $\varepsilon/2$ each.
Use $|x_m-x_n|\le|x_m-L|+|x_n-L|$ to obtain a Cauchy estimate.
For a direct infimum cut, take $L$ as all lower bounds and
$U=\R\setminus L$; prove separation, then use
$c\le s+\varepsilon$ for every $\varepsilon>0$ to obtain $c\le s$.
Do not invoke LUBP if the question forbids it.

\formulagroup{Missing hypotheses and missing limits}
Use $(0,1/n)$ to drop closedness, $[n,\infty)$ to drop boundedness,
and shrinking rational-ended intervals about an irrational to expose
incompleteness of $\Q$.
Rational decimals $q_n$ with $|q_n-\alpha|<10^{-n}$ are Cauchy
but have no rational limit when $\alpha$ is irrational.
Small covering length alone does not prove incompleteness.

\formulagroup{Root uniqueness and equality conditions}
For nonnegative expressions, raise to a common integer power and use
uniqueness; for rational exponents use a common denominator.
Even powers do not identify signed reals.
$ (\sqrt a-\sqrt b)^2=a+b-2\sqrt{ab}\ge0$ proves AM--GM,
with equality exactly at $a=b$.
For reverse triangle, write $x=(x-y)+y$, then swap $x,y$;
equality holds exactly when $xy\ge0$.

\ifdefined\standaloneformulasheet
\sheetnewtonandquadratic
\fi
\formulagroup{Recurrences and fixed points}
Simplify sums first. Prove an invariant interval by induction and
factor $g(x)-x$ to get monotonicity. Bounded monotonicity proves
convergence; only then solve $L=g(L)$, retaining roots compatible
with the bounds. Positive terms alone do not imply $L>0$.
For $g(x)=(\tfrac32x^2+2x)/(1+x^2)$,
\[
  2-g(x)=\frac{(x-2)^2}{2(1+x^2)},\qquad
  g(x)-x=\frac{x(2-x)(2x+1)}{2(1+x^2)}.
\]
The general Newton map for $x-a-1/x$ is
$(ax^2+2x)/(1+x^2)$; its positive limiting fixed point is
$(a+\sqrt{a^2+4})/2$.
Fibonacci ratios satisfy $r_n=1+1/r_{n-1}$ for $n\ge3$
and $1\le r_n\le2$ for $n\ge2$;
if convergent, their limit is $(1+\sqrt5)/2$.

\formulagroup{Integer remainder, oscillation, and $n$th roots}
If $e=p/q$, take $N\ge\max\{q,2\}$; $N!(e-s_N)$ is an integer in $(0,1)$.
For periodic behaviour, split indices modulo the period; an infinite
subsequence contains infinitely many indices in one class.
For $n$th roots, factor the dominant term and squeeze using
$\max\{u,v\}\le u+v\le2\max\{u,v\}$ and $c^{1/n},n^{1/n}\to1$.

\formulagroup{Finite tails, strict gaps, and last checks}
Telescope finite partial sums and keep the last term.
For alternating tails, pair in both parities, retaining any unpaired
term, to get the first-omitted-term bound.
Rationalise arbitrary real shifts rather than assuming cancellation.
Choose $\limsup t_n<r<1$ for an eventual upper bound or
$1<r<\liminf t_n$ for a lower bound.
Test value $1$ is inconclusive; test power-series endpoints separately.
Finite prefixes do not change convergence, but can change extrema and sums.
